\chapter{The Formatter and Doc Generator}\label{ch:formatter-docgen}

\index{formatter}\index{clat fmt}\index{doc generator}\index{clat doc}

Every codebase drifts.
One developer uses tabs, another uses spaces.
Someone wraps function arguments at 80 columns; someone else at 120.
Doc comments get written in three different styles.
A week later, the ``git blame'' for every file reads like a United Nations assembly.
Lattice solves this with two built-in tools: \lstinline|clat fmt| for formatting source code and \lstinline|clat doc| for generating documentation.
Both are first-class citizens of the \lstinline|clat| CLI---no third-party installation required.

%% ─────────────────────────────────────────────────────────────
\section{\texttt{clat fmt} --- Formatting Your Code}\label{sec:clat-fmt}
\index{clat fmt}

Let's start with a messy file.
Suppose a colleague has handed us \filename{sensor.lat} with inconsistent spacing, missing indentation, and some creative use of whitespace:

\begin{lstlisting}[language=Lattice, caption={Before formatting: messy whitespace and indentation}]
fn read_temperature(  sensor_id:Int,unit: String   ) -> Float{
    let raw=sensor_id *   0.48+2.1
if unit ==  "celsius"{
return raw
}else{
    return raw*9.0/5.0+32.0
}}
\end{lstlisting}

One command brings order to the chaos:

\begin{lstlisting}[language=Lattice, caption={Terminal: running clat fmt}]
// In your terminal:
// $ clat fmt sensor.lat
\end{lstlisting}

After formatting, the file is rewritten in place:

\begin{lstlisting}[language=Lattice, caption={After formatting: clean and consistent}]
fn read_temperature(sensor_id: Int, unit: String) -> Float {
    let raw = sensor_id * 0.48 + 2.1
    if unit == "celsius" {
        return raw
    } else {
        return raw * 9.0 / 5.0 + 32.0
    }
}
\end{lstlisting}

Notice what the formatter did:

\begin{itemize}
  \item Normalized spacing around operators (\lstinline|=|, \lstinline|*|, \lstinline|+|) to exactly one space on each side.
  \item Enforced the \emph{attached brace style}: the opening \lstinline|\{| appears on the same line as the \lstinline|fn|, \lstinline|if|, and \lstinline|else| keywords, preceded by a single space.
  \item Inserted a space after the colon in type annotations (\lstinline|sensor_id: Int|) and after commas.
  \item Applied consistent four-space indentation at each nesting level.
  \item Placed the \lstinline|else| clause on the same line as the closing brace of the \lstinline|if| block (\lstinline|} else {|).
  \item Ensured a single trailing newline at the end of the file.
\end{itemize}

\begin{defnbox}{clat fmt}
The \lstinline|clat fmt| command reads a Lattice source file, reformats it using the canonical style, and writes the result back. It operates on a character-by-character scan of the source (implemented in \filename{src/formatter.c}), preserving all comments and string contents while normalizing whitespace and indentation. The default indentation width is four spaces, and the default target line width is 100 columns.
\end{defnbox}

\subsection{Formatting Multiple Files}\label{sec:fmt-multiple-files}

You can format an entire project in one go by passing multiple files:

\begin{lstlisting}[language=Lattice, caption={Formatting several files at once}]
// $ clat fmt src/main.lat src/utils.lat lib/math.lat
\end{lstlisting}

\subsection{Check Mode: CI-Friendly Verification}\label{sec:fmt-check}
\index{clat fmt!check mode}

In a continuous integration pipeline, you don't want to \emph{rewrite} files---you want to \emph{verify} that they're already formatted. The \lstinline|--check| flag does exactly this:

\begin{lstlisting}[language=Lattice, caption={Using {-}{-}check in CI}]
// $ clat fmt --check src/main.lat
// If already formatted: exits with code 0
// If not: exits with code 1 and prints a diff
\end{lstlisting}

Under the hood, \lstinline|--check| calls the \lstinline|lat_format_check()| function (see \filename{include/formatter.h}), which formats the source and compares the result to the original. If they match, the file is already canonical.

\begin{tipbox}{Add Formatting to Your CI Pipeline}
A common pattern for Lattice projects:
\begin{lstlisting}[language=Lattice]
// In your CI configuration:
// $ clat fmt --check src/ lib/
\end{lstlisting}
This ensures every pull request follows the same style. No more bikeshedding in code reviews about where the braces go.
\end{tipbox}

\subsection{Reading from Standard Input}\label{sec:fmt-stdin}
\index{clat fmt!stdin}

The formatter can also read from standard input with the \lstinline|--stdin| flag.
This is useful for editor integrations and piped workflows:

\begin{lstlisting}[language=Lattice, caption={Formatting via stdin}]
// $ cat messy.lat | clat fmt --stdin
// Formatted output goes to stdout
\end{lstlisting}

The \lstinline|lat_format_stdin()| function in \filename{src/formatter.c} reads all of standard input into a buffer, formats it, and writes the result to standard output.

%% ─────────────────────────────────────────────────────────────
\section{\texttt{--width} for Configurable Line Width}\label{sec:fmt-width}
\index{clat fmt!width}\index{line width}

The default line width target is 100 columns---a comfortable middle ground for modern screens.
But preferences vary.
Some teams work on laptops with split panes and prefer 80 columns.
Others have ultrawide monitors and 120 feels right.

The \lstinline|--width| flag lets you choose:

\begin{lstlisting}[language=Lattice, caption={Setting line width to 80 columns}]
// $ clat fmt --width 80 src/main.lat
\end{lstlisting}

\begin{lstlisting}[language=Lattice, caption={Setting line width to 120 columns}]
// $ clat fmt --width 120 src/main.lat
\end{lstlisting}

When the formatter encounters a line that exceeds the target width, it makes wrapping decisions---for instance, breaking long function parameter lists across multiple lines.
Passing zero (or omitting the flag) uses the default of 100, as defined by the \lstinline|DEFAULT_LINE_WIDTH| constant in \filename{src/formatter.c}.

\begin{notebox}{Width is a Target, Not a Hard Limit}
The \lstinline|--width| value is a \emph{target}, not an absolute ceiling.
Some constructs---like a long string literal or a deeply nested expression---cannot be broken further without changing program semantics.
In those cases, the formatter will exceed the target rather than produce incorrect code.
\end{notebox}

Let's see the width flag in action.
Consider a function with many parameters:

\begin{lstlisting}[language=Lattice, caption={A function with many parameters}]
fn create_report(title: String, author: String, created_at: String, tags: Array, body: String, draft: Bool) -> Map {
    let report = Map::new()
    report["title"] = title
    report["author"] = author
    report["created_at"] = created_at
    report["tags"] = tags
    report["body"] = body
    report["draft"] = draft
    return report
}
\end{lstlisting}

At \lstinline|--width 80|, the formatter may wrap the parameter list to keep lines within bounds.
At \lstinline|--width 120|, the entire signature might fit on a single line.

\begin{tipbox}{Team Convention}
Pick a width for your project and stick with it.
Add it to your project's build script or Makefile:
\begin{lstlisting}[language=Lattice]
// Makefile target:
// fmt:
//     clat fmt --width 100 src/ lib/
\end{lstlisting}
\end{tipbox}

%% ─────────────────────────────────────────────────────────────
\section{\texttt{///} Doc Comments}\label{sec:doc-comments}
\index{doc comments}\index{triple-slash comments}

Lattice uses triple-slash (\lstinline|///|) comments for documentation.
Unlike regular \lstinline|//| comments, doc comments are extracted by the documentation generator and attached to the declaration that follows them.

\begin{lstlisting}[language=Lattice, caption={Doc comments on a function}]
/// Calculate the body mass index for a patient.
/// Returns the BMI as a Float value.
fn calculate_bmi(weight_kg: Float, height_m: Float) -> Float {
    return weight_kg / (height_m * height_m)
}
\end{lstlisting}

Doc comments can span multiple consecutive lines.
The \lstinline|///| prefix and one optional space after it are stripped, and the remaining text becomes the documentation body.

\subsection{Documenting All Declaration Types}\label{sec:doc-comments-types}

Doc comments work on every kind of top-level declaration: functions, structs, enums, traits, impl blocks, and even variables.

\begin{lstlisting}[language=Lattice, caption={Doc comments on structs and enums}]
/// A patient record in the healthcare system.
/// Contains demographic and medical information.
struct Patient {
    /// The patient's full legal name.
    name: String,
    /// Age in years.
    age: Int,
    /// Blood type (e.g., "A+", "O-").
    blood_type: String
}

/// Possible outcomes of a medical test.
enum TestResult {
    /// The test came back normal.
    Negative,
    /// The test detected the condition.
    Positive,
    /// The sample was insufficient or contaminated.
    Inconclusive
}
\end{lstlisting}

Notice that each struct field and enum variant can have its own doc comment.
The documentation generator (covered in the next section) extracts these and associates them with the correct field or variant.

\subsection{Module-Level Doc Comments}\label{sec:module-doc}
\index{doc comments!module-level}

If a \lstinline|///| comment block appears at the very top of a file---before any declaration---it is treated as a \emph{module-level} doc comment.
This is the place to describe what the file or module does:

\begin{lstlisting}[language=Lattice, caption={Module-level documentation}]
/// The geometry module provides functions for calculating
/// areas and perimeters of common shapes.
///
/// All functions accept dimensions in the same unit
/// and return results in squared units where applicable.

fn circle_area(radius: Float) -> Float {
    return 3.14159265 * radius * radius
}

fn rectangle_area(width: Float, height: Float) -> Float {
    return width * height
}
\end{lstlisting}

The doc extractor in \filename{src/doc\_gen.c} distinguishes module-level docs from declaration docs by checking whether the doc block is followed by a declaration keyword (\lstinline|fn|, \lstinline|struct|, \lstinline|enum|, \lstinline|trait|, \lstinline|impl|, \lstinline|flux|, \lstinline|fix|, \lstinline|let|, or \lstinline|export|).
If not, it's treated as the module description.

\subsection{Doc Comment Conventions}\label{sec:doc-conventions}

While Lattice doesn't enforce a specific doc comment format, the community has settled on a few conventions:

\begin{lstlisting}[language=Lattice, caption={Recommended doc comment style}]
/// Open a database connection to the given host.
///
/// The connection uses a pooled TCP socket under the hood.
/// If the host is unreachable, returns an error after
/// the specified timeout.
///
/// Parameters are validated before connecting. An empty
/// host string will produce an immediate error.
fn connect(host: String, port: Int, timeout_ms: Int) -> Map {
    // ... implementation
    return Map::new()
}
\end{lstlisting}

\begin{itemize}
  \item First line: a brief one-sentence summary.
  \item Blank line (an empty \lstinline|///| line) separating the summary from the longer description.
  \item Additional paragraphs for context, usage notes, or caveats.
\end{itemize}

\begin{warnbox}{Regular Comments Are Ignored}
Only \lstinline|///| comments are picked up by \lstinline|clat doc|.
Regular \lstinline|//| comments and block comments (\lstinline|/* ... */|) are treated as implementation notes and excluded from generated documentation.
If you write documentation with \lstinline|//| instead of \lstinline|///|, it won't appear in the output.
\end{warnbox}

%% ─────────────────────────────────────────────────────────────
\section{\texttt{clat doc} --- Generating Documentation}\label{sec:clat-doc}
\index{clat doc}\index{documentation generation}

The \lstinline|clat doc| command extracts doc comments from one or more \lstinline|.lat| files and renders them as human-readable documentation.

\begin{lstlisting}[language=Lattice, caption={Basic usage of clat doc}]
// $ clat doc geometry.lat
\end{lstlisting}

By default, this prints Markdown to standard output.
Let's walk through a complete example.
Suppose we have \filename{math\_utils.lat}:

\begin{lstlisting}[language=Lattice, caption={A well-documented utility module}]
/// Utility functions for common mathematical operations.

/// Clamp a value to the range [low, high].
/// If the value is below low, returns low.
/// If above high, returns high.
fn clamp(value: Float, low: Float, high: Float) -> Float {
    if value < low {
        return low
    }
    if value > high {
        return high
    }
    return value
}

/// Linear interpolation between two values.
/// When t is 0.0, returns a. When t is 1.0, returns b.
fn lerp(a: Float, b: Float, t: Float) -> Float {
    return a + (b - a) * t
}

/// A 2D point in Cartesian space.
struct Point {
    /// Horizontal coordinate.
    x: Float,
    /// Vertical coordinate.
    y: Float
}

/// Compute the Euclidean distance between two points.
fn distance(p1: Point, p2: Point) -> Float {
    let dx = p2.x - p1.x
    let dy = p2.y - p1.y
    return (dx * dx + dy * dy) ** 0.5
}
\end{lstlisting}

Running \lstinline|clat doc math_utils.lat| produces:

\begin{lstlisting}[language=Lattice, caption={Generated Markdown output (abbreviated)}]
// Utility functions for common mathematical operations.
//
// ## Functions
//
// ### `clamp(value: Float, low: Float, high: Float) -> Float`
//
// Clamp a value to the range [low, high].
// If the value is below low, returns low.
// If above high, returns high.
//
// ### `lerp(a: Float, b: Float, t: Float) -> Float`
//
// Linear interpolation between two values.
// When t is 0.0, returns a. When t is 1.0, returns b.
//
// ## Structs
//
// ### `struct Point`
//
// A 2D point in Cartesian space.
//
// | Field | Type    | Description            |
// |-------|---------|------------------------|
// | `x`   | `Float` | Horizontal coordinate. |
// | `y`   | `Float` | Vertical coordinate.   |
\end{lstlisting}

The documentation generator groups items by kind: functions, structs, enums, traits, implementations, and variables.
Each section includes the full signature and associated documentation.

\subsection{Documenting a Directory}\label{sec:doc-directory}

Pass a directory to document all \lstinline|.lat| files inside it:

\begin{lstlisting}[language=Lattice, caption={Documenting an entire directory}]
// $ clat doc src/
\end{lstlisting}

When processing multiple files, each file gets its own heading in the output. This makes it straightforward to generate documentation for an entire project.

\subsection{Writing Output to Files}\label{sec:doc-output-dir}
\index{clat doc!output directory}

For larger projects, you'll want to write documentation to files rather than standard output.
The \lstinline|-o| (or \lstinline|--output|) flag specifies an output directory:

\begin{lstlisting}[language=Lattice, caption={Writing docs to a directory}]
// $ clat doc -o docs/ src/
// wrote docs/math_utils.md
// wrote docs/geometry.md
// wrote docs/string_helpers.md
\end{lstlisting}

Each input file produces a corresponding output file.
The \lstinline|.lat| extension is replaced with the appropriate extension for the chosen format (\lstinline|.md| for Markdown, \lstinline|.json| for JSON, \lstinline|.html| for HTML, or \lstinline|.txt| for plain text).
The output directory is created automatically if it doesn't exist.

%% ─────────────────────────────────────────────────────────────
\section{Output Formats: Markdown, JSON, HTML}\label{sec:doc-formats}
\index{clat doc!output formats}

The documentation generator supports four output formats, each suited to a different use case.

\subsection{Markdown (Default)}\label{sec:doc-markdown}
\index{clat doc!Markdown}

Markdown is the default because it's the lingua franca of technical documentation---readable as plain text, renderable on GitHub, GitLab, and most wikis.

\begin{lstlisting}[language=Lattice, caption={Generating Markdown output}]
// $ clat doc --md src/server.lat
// (or simply: clat doc src/server.lat)
\end{lstlisting}

The Markdown renderer produces headings, code-formatted signatures, tables for struct fields, and bulleted lists for enum variants.

\subsection{Plain Text}\label{sec:doc-text}
\index{clat doc!plain text}

For terminal viewing or piping to other tools:

\begin{lstlisting}[language=Lattice, caption={Generating plain text output}]
// $ clat doc --text src/server.lat
\end{lstlisting}

The text format uses simple labels like \texttt{FUNCTIONS}, \texttt{STRUCTS}, and \texttt{ENUMS}, with dashed underlines and indented descriptions.

\subsection{JSON}\label{sec:doc-json}
\index{clat doc!JSON}

The JSON format is designed for machine consumption.
It's what you'd use to build a custom documentation website, IDE plugin, or search index:

\begin{lstlisting}[language=Lattice, caption={Generating JSON output}]
// $ clat doc --json src/math_utils.lat
\end{lstlisting}

The output is a JSON array where each element represents a source file.
Each file object contains \lstinline|"file"|, an optional \lstinline|"module_doc"|, and an \lstinline|"items"| array.
Each item has a \lstinline|"kind"| (\lstinline|"function"|, \lstinline|"struct"|, \lstinline|"enum"|, \lstinline|"trait"|, \lstinline|"impl"|, or \lstinline|"variable"|), a \lstinline|"name"|, a \lstinline|"line"| number, and optional \lstinline|"doc"| text, plus kind-specific fields like \lstinline|"params"| for functions or \lstinline|"fields"| for structs.

\begin{lstlisting}[language=Lattice, caption={Example JSON output (abbreviated)}]
// [
//   {
//     "file": "math_utils.lat",
//     "module_doc": "Utility functions for ...",
//     "items": [
//       {
//         "kind": "function",
//         "name": "clamp",
//         "line": 4,
//         "doc": "Clamp a value to the range ...",
//         "params": [
//           {"name": "value", "type": "Float"},
//           {"name": "low", "type": "Float"},
//           {"name": "high", "type": "Float"}
//         ],
//         "return_type": "Float"
//       }
//     ]
//   }
// ]
\end{lstlisting}

\begin{tipbox}{Building a Documentation Website}
The JSON output pairs well with static site generators.
Extract docs with \lstinline|clat doc --json src/ > api.json|, then feed \lstinline|api.json| into your favorite template engine to produce a searchable HTML site.
\end{tipbox}

\subsection{HTML}\label{sec:doc-html}
\index{clat doc!HTML}

For a self-contained, styled documentation page:

\begin{lstlisting}[language=Lattice, caption={Generating HTML documentation}]
// $ clat doc --html -o docs/ src/
\end{lstlisting}

The HTML renderer produces a complete web page with a dark theme, syntax-highlighted signatures, and interactive hover effects.
Function signatures are color-coded---keywords in purple, types in gold, function names in blue---making it easy to scan a large API surface.

Struct fields render as tables, enum variants as styled lists, and trait methods as indented signatures under their parent trait.

\subsection{The \texttt{--doc-format} Flag}\label{sec:doc-format-flag}

As an alternative to the shorthand flags, you can use \lstinline|--doc-format| with an explicit format name:

\begin{lstlisting}[language=Lattice, caption={Using {-}{-}doc-format}]
// $ clat doc --doc-format html src/
// $ clat doc --doc-format json src/
// $ clat doc --doc-format md src/
// $ clat doc --doc-format text src/
\end{lstlisting}

This is useful in scripts where you want to parameterize the format.

\subsection{How Extraction Works Under the Hood}\label{sec:doc-extraction-internals}

The documentation extractor in \filename{src/doc\_gen.c} uses its own lightweight scanner (\lstinline|DocScanner|) rather than the full Lattice lexer.
This is a deliberate design choice: the scanner needs to preserve doc comments while skipping implementation details, so it only recognizes enough syntax to identify declarations.

The process works as follows:

\begin{enumerate}
  \item The scanner reads through the source file, skipping whitespace and regular comments.
  \item When it encounters a \lstinline|///| line, it reads consecutive doc comment lines into a single block, stripping the prefix and one optional leading space.
  \item After the doc block, the scanner looks for a declaration keyword (\lstinline|fn|, \lstinline|struct|, \lstinline|enum|, \lstinline|trait|, \lstinline|impl|, or a phase keyword).
  \item For each declaration, the scanner extracts the name and kind-specific details: function parameters and return types, struct fields with their types, enum variants with optional tuple parameters, trait method signatures, and impl block methods.
  \item The result is a \lstinline|DocFile| structure (defined in \filename{include/doc\_gen.h}) containing a list of \lstinline|DocItem| values that the renderers consume.
\end{enumerate}

Items that lack doc comments are still extracted for functions, structs, enums, and traits.
Variables, however, are only included if they have a doc comment, since most local or temporary variables aren't part of a module's public API.

\begin{notebox}{Test Blocks Are Skipped}
The doc generator recognizes \lstinline|test "name" \{ ... \}| blocks and skips them entirely.
Tests are implementation details, not part of your public documentation.
\end{notebox}

%% ─────────────────────────────────────────────────────────────
\section*{Exercises}\label{sec:formatter-docgen-exercises}
\addcontentsline{toc}{section}{Exercises}

\begin{enumerate}
  \item Take one of the Lattice files you've written in earlier chapters and run \lstinline|clat fmt| on it. Compare the before and after. Were there any formatting choices that surprised you? Try running with \lstinline|--width 80| and \lstinline|--width 120| to see how the output changes.

  \item Create a small library module (perhaps a \lstinline|temperature.lat| with conversion functions) and document every function, struct, and constant with \lstinline|///| comments. Then run \lstinline|clat doc --html -o docs/ temperature.lat| and open the resulting HTML in a browser.

  \item Write a shell script (or Makefile target) that runs \lstinline|clat fmt --check| on all \lstinline|.lat| files in a directory tree and exits with a non-zero code if any file needs formatting. This is the skeleton of a CI formatting check.

  \item Generate JSON documentation for a multi-file project and write a small Lattice script that reads the JSON and counts the total number of documented functions, structs, and enums across all files.
\end{enumerate}

%% ─────────────────────────────────────────────────────────────
\section*{What's Next}\label{sec:formatter-docgen-next}
\addcontentsline{toc}{section}{What's Next}

With clean formatting and thorough documentation in hand, the natural next step is to make sure everything actually \emph{works}.
In \cref{ch:testing}, we'll explore Lattice's built-in testing system---inline \lstinline|test| blocks that live right alongside your code, a rich set of assertion functions, and the \lstinline|clat test| command that discovers and runs your entire test suite.
